\section{Conclusion}
\label{sec:conclusion}

This paper introduced transition certificates for verifying Linear Temporal Logic (LTL) properties of discrete-time dynamical systems with continuous state spaces. The method separates the proof obligation into transition safety certificates, which make selected automaton transitions unreachable, and transition persistence certificates, which ensure that accepting transitions cannot occur infinitely often. This separation reduces conservatism relative to state-triplet certificates on the reported persistence-style cases and gives a smaller synthesis domain than closure certificates for the same automaton product.

The experiments show where this proof structure is useful. On the reported Kuramoto oscillator and two-room temperature benchmarks, transition certificates obtain verified certificates within the bounded synthesis budget in cases where the compared closure-certificate encoding is more expensive or times out. These results support the claimed proof and synthesis advantages for the tested specifications. The claim is limited to the reported formulas, systems, encodings, and solver budget.

Several limitations remain. Neural-network certificate templates still depend on SMT solving and can become expensive for high-dimensional or highly nonlinear systems. The construction of the nondeterministic Buechi automaton also inherits the usual exponential worst-case dependence on the LTL formula. Extending the certificate conditions to continuous-time systems and connecting them to controller synthesis are natural next steps.
